\section{Discussion}
\label{sec:discussion}

\paragraph{What a responder can take from this.} Three findings are
actionable without any further work. A robustness claim for an aerial
person detector is not portable across altitude, so a number quoted from
one capture regime should not be used to select a system for another.
Low light is not a graded degradation for these detectors but a cliff,
and at roughly 4.5 stops of exposure loss the systems return nothing at
all. And because the failure is a loss of detections rather than
mislocalisation, an empty feed under degraded conditions carries no
information: it is indistinguishable from a correctly searched empty
frame, so it should not be treated as coverage.

\paragraph{Why per-condition reporting is the unit.} Every headline in
this paper is invisible under averaging. The span in
Table~\ref{tab:spectrum} runs from a 3.6-point loss to a total one
inside one model and one severity. The rain reversal between datasets
cancels if the two are pooled. The mitigation trade-off in
Table~\ref{tab:tradeoff} changes sign across architectures. We therefore
release results per condition, per model and per dataset with intervals,
and would encourage the same shape for any robustness result intended to
inform a deployment decision.

\paragraph{Limitations.} Four bound what is claimed here. First, the
corruptions are synthetic, applied to real imagery, following the
ImageNet-C and Foggy Cityscapes
lineage~\cite{hendrycks2019imagenetc,sakaridis2018foggy}; this is what
makes controlled per-condition attribution possible, since disaster
imagery rarely has a matched clean counterpart of the same scene, but it
is also the ceiling of the method. Second, the VisDrone spot check is
40 images and is suggestive rather than a substitute for systematic
validation against real adverse-condition search-and-rescue footage,
which is made harder by how rarely such footage is released: aerial
disaster imagery exists for related tasks such as fire segmentation but
not with person annotations, itself a gap this release may help motivate
closing. Third, we evaluate two datasets and three detectors,
prioritising depth of protocol over breadth of coverage, and the
multi-seed confirmation covers one model on one dataset. Fourth,
corruptions are applied one at a time, whereas a real night flood
presents several at once; the released engine composes them, but we do
not report the combinatorial grid and do not assume the effects are
additive.

\paragraph{Ethics and responsible release.} We use only publicly
available research datasets under their own terms and collect no
human-subjects data. We add no imagery and no annotations of people, and
annotate no identity, biometric or demographic attribute; persons appear
only as generic boxes inherited from the sources. Automated aerial
person detection is dual-use, but this work introduces no detector, and
its finding, that existing detectors fail under degraded conditions,
extends no surveillance capability. The clearest risk in publishing a
negative result about a safety-of-life system is that it reads as a
reason to abandon the technology. It is not: in clear conditions these
systems provide coverage ground teams cannot match. The larger risk is
the status quo, in which robustness is reported only under clear
conditions and an operator has no basis for knowing when the feed has
stopped being informative.

\section{Released Artefacts}
\label{sec:release}

The release is organised so that a reader can check any number in this
paper against its source without rerunning anything, and rerun anything
they do not trust.

\begin{itemize}\itemsep1pt \parskip0pt \topsep2pt
\item \textbf{Corruption engine and parameters.} All nine corruptions,
their per-severity parameters, the declared calibration statistic for
each, and the calibration harness that verifies monotonicity. Seeding is
a stable hash of image identity, corruption, severity and global seed,
so the corrupted benchmark regenerates bit-identically from the source
datasets on any platform. No imagery is redistributed.
\item \textbf{Evaluation pipeline.} Tiling, prediction remapping,
class-agnostic NMS, the shared IoU-matching evaluator, the
F1-maximising threshold selection, and the bootstrap interval
computation, as run.
\item \textbf{Per-condition results with intervals.} One row per
(model, dataset, corruption, severity) for all 168 conditions, carrying
recall with its bootstrap 95\% interval, in a single flat CSV. The
mitigation and multi-seed arms ship as separate tables in the same
column shape, so they can be joined on the condition key.
\item \textbf{Raw sweep records.} The full per-run table underlying the
above, including precision, F1, mAP$_{50}$, mAP$_{50:95}$, recall by
COCO area band, ground-truth and true-positive counts, the frozen
confidence threshold actually used, and the git commit and config hash
for each run.
\item \textbf{Training configurations.} Every config referenced here,
including the three mitigation arms and the two additional seeds, at the
exact settings that produced the reported numbers.
\end{itemize}

Every table in this paper is generated from those files by a checked
script rather than transcribed, and the build fails if any printed value
disagrees with its source record.

\section{Conclusion}

AegisBench measures what aerial search-and-rescue person detectors do
under the conditions that cause them to be launched, using physically
modelled corruptions at calibrated severities and operating points a
fielded system could actually have chosen in advance. Degradation across
168 conditions is uneven, does not transfer between capture regimes, and
is not summarised by any single number. Low light drives all three
architectures to exactly zero recall on both datasets, a result that
survives a threshold-independent check, a three-seed repeat and a test
on real night imagery, and the failure takes the form of detections
disappearing rather than boxes drifting, which is the form least visible
to an operator. Targeted augmentation moves that floor at a severity it
never trained on, confirming the deficit is one of data coverage rather
than architecture, but its cost on untrained conditions changes sign
across architectures. The measurement, not the mitigation, is the
contribution: until robustness for this task is reported per condition,
a deployment decision is being made on a number that does not describe
the mission.
